% %
% LAYOUT_E.TEX - Short description of REFMAN.CLS
%                                       99-03-20
%
%  Updated for REFMAN.CLS (LaTeX2e)
%
\documentclass[twoside,a4paper]{refart}
\usepackage{makeidx}
\usepackage{ifthen}
\usepackage{comment}
\usepackage{amsmath}
\usepackage{graphicx}
\usepackage{url}
\usepackage{lscape}

% ifthen wird vom Bild von N.Beebe gebraucht!

\def\bs{\char'134 } % backslash in \tt font.
\newcommand{\ie}{i.\,e.,}
\newcommand{\eg}{e.\,g..}
\renewcommand{\mod}{\text{\textbf{\textit{Complex Calc}}}}

\DeclareRobustCommand\cs[1]{\texttt{\char`\\#1}}

% Configuración para el título personalizado CORREGIDA
% Título con imagen - FORMA CORRECTA
\title{%
    \begin{minipage}{0.9\textwidth}
        \raggedright
        \LARGE\bfseries
        $\mod$ (Calculadora de Sistemas de Ecuaciones Lineales Complejos)
    \end{minipage}%
    \hfill%
    \begin{minipage}{0.25\textwidth}
        \centering
        
\includegraphics[width=0.7\linewidth]{Projects/Complex calc/IE.png} % Cambia por tu imagen
    \end{minipage}
}



\author{
Contreras Reyes Orlando \\
Garcia Barrera Iker Emiliano \\
Lara Hernández Kevin Adrián \\
López Rodríguez Roberto
}
\date{
Universidad Autónoma de Aguascalientes \\
Centro de Ciencias Básicas \\
Curso: Circuitos Eléctricos en CA \\
Fecha de publicación: Noviembre 2025 \\
Versión: 2.7
}

\emergencystretch1em  %

\pagestyle{myfootings}
\markboth{$\mod$}% Foot page
         {$\mod$} 

\makeindex 

\setcounter{tocdepth}{2}

\begin{document}
\maketitle

\begin{comment}
    2 Elementos de la Interfaz
    3 Modos de Uso
    3.1 Complejos Polar
    3.2 Complejos Rectangular
    3.3 Reales
    6. Licencia de Uso
\end{comment}

\tableofcontents

\newpage


%%%%%%%%%%%%%%%%%%%%%%%%%%%%%%%%%%%%%%%%%%%%%%%%%%%%%%%%%%%%%%%%%%%%
\section{Introducción}

\subsection{¿Qué es?}
Una calculadora es una herramienta sumamente útil para distintos ámbitos de la ciencia y la ingeniería, ya que facilita la realización de operaciones que normalmente tomarían mucho tiempo al realizarlas manualmente. En esta ocasión se presenta a los usuarios una opción sencilla y rápida de utilizar para resolver problemas que implican el uso de números complejos y sistemas de ecuaciones.

\textbf{Complex Calc} es una gran opción debido a su facilidad de uso, su capacidad de automatizar cálculos y su orientación al análisis en circuitos eléctricos. El propósito principal de esta herramienta es ofrecer una solución eficiente que reduzca el tiempo de trabajo, minimice errores humanos y permita obtener resultados inmediatos. Además, esta calculadora está pensada tanto para estudiantes como para profesionales que requieren una solución directa para problemas basados en operaciones complejas, sin la necesidad de recurrir a procesos extensos o software especializado de mayor complejidad.

\subsection{Números Complejos}
Los números complejos permiten abordar problemas desde otra perspectiva matemática, proporcionando una forma de representar cantidades que involucran magnitudes y fases. Un ejemplo común de su aplicación es el análisis de señales en corriente alterna (CA). Para su estudio existen diversos métodos, entre los más populares se encuentran los métodos basados en las Leyes de Kirchhoff: Nodos y Mallas.

Estos métodos requieren operaciones matemáticas tales como sistemas de ecuaciones, sumas, restas, multiplicaciones y divisiones de fasores o números complejos. Debido a esto, resulta útil disponer de una herramienta que permita encontrar soluciones de manera rápida y efectiva. La calculadora desarrollada en este proyecto, \textbf{Complex Calc}, tiene como finalidad facilitar el análisis de circuitos eléctricos mediante el uso de procedimientos automatizados que simplifican el cálculo de ecuaciones en el dominio de los números complejos.

\subsection{Aplicaciones Prácticas}
El uso de números complejos y fasores no se limita únicamente al análisis de circuitos. Entre las aplicaciones más relevantes se encuentran:

\begin{itemize}
\item Análisis de impedancias en sistemas trifásicos.
\item Simulación de sistemas de potencia.
\item Estudio de filtros y respuesta en frecuencia.
\item Control de máquinas eléctricas rotativas.
\item Comunicación y transmisión de señales.
\end{itemize}

En todos estos casos, la interpretación correcta de los resultados depende en gran medida del manejo adecuado de los números complejos. Una herramienta como \textbf{Complex Calc} facilita la toma de decisiones, el diseño de circuitos y el análisis de resultados experimentales o teóricos.

\newpage
\section{Elementos de La interfaz}
La interfaz consta de una pantalla con un historial y diferentes botones que permiten un manejo eficiente de $\mod$ como se puede apreciar en la figura \ref{fig:UI}
\begin{figure}[!h]
    \centering
    \includegraphics[width=1.30\linewidth]{Projects/Complex calc/GUI.png}
    \caption{Interfaz de $\mod$}
    \label{fig:UI}
\end{figure}

\section{Modos de uso}
Como demostración del uso de la calculadora se resolvera el siguiente circuito [Fig \ref{fig:Mesh}] sugerido por el \textbf{Doctor Luis Alejandro Flores Oropeza}. Se realizará las ecuaciones de forma polar y rectangular demostrando la versatilidad de $\mod$.
\begin{figure}[!h]
    \centering
    \includegraphics[width=0.8\linewidth]{Projects/Mallas2.png}
    \caption{Circuito AC 2 Mallas}
    \label{fig:Mesh}
\end{figure}
\subsection{Sistemas de Ecuaciones Complejos Rectangulares}
Considerando el circuito [Fig \ref{fig:Mesh}] se llega a la ecuación
$$
\begin{cases}
\begin{aligned}
(15 + j5)\,I_1 &\;-\; 15\,I_2 &=&\; 10 + j0 \\[6pt]
-15\,I_1 &\;+\; (15 - j10)\,I_2 &=&\; 0 - j5
\end{aligned}
\end{cases}
\Rightarrow
\left[
\begin{array}{cc|c}
15 + j5 & -15      & 10 + j0 \\[6pt]
-15     & 15 - j10 & 0 - j5
\end{array}
\right]
$$


Para ello seleccionaremos el botón de cambiar tamaño y abrirá un menú [Fig \ref{fig:tamano}] donde escribiremos el valor de 2 siendo nuestra ecuación de 2 incógnitas y procederemos a darle en el boton ok
\begin{figure}[h]
    \centering
    \includegraphics[width=0.4\linewidth]{Projects//Complex calc/size.png}
    \caption{Cambiar tamaño}
    \label{fig:tamano}
\end{figure}


Esto abrirá un nuevo menu con un total de 6 espacios disponibles, para ello podemos seleccionar casilla por casilla o usar la tecla $<\text{TAB}>$ para navegar entre casillas e ingresar nuestro sistema de ecuaciones.
Finalmente daremos en el boton resolver para observar los valores de las corrientes de el circuito \ref{fig:Mesh} en la figura \ref{fig:SolRect}


\begin{figure}[h]
        \centering
        \includegraphics[width=0.9\linewidth]{Projects//Complex calc/RectSol.png}
        \caption{Solución en forma rectangular}
        \label{fig:SolRect}
\end{figure}
\newpage

\subsection{Sistemas de Ecuaciones Complejos Polares}
Para la sección de polares haremos el mismo circuito de la figura [Fig \ref{fig:Mesh}] pero con fasores/complejos polares, para ello tendremos que convertir los valores de rectangular a su forma polar.
\[
\begin{aligned}
15 + j5 &= 15.81\angle 18.43^\circ \\[4pt]
-15 &= 15\angle 180^\circ \\[4pt]
15 - j10 &= 18.03\angle -33.69^\circ \\[4pt]
10\angle 0^\circ &\text{ \; (ya está en forma polar)} \\[4pt]
-5\angle 90^\circ &= 5\angle -90^\circ
\end{aligned}
\]
Con esto llegamos a una sistema de ecuaciones de la siguiente forma.
$$
\left[
\begin{array}{cc|c}
15.81\angle 18.43^\circ & 15\angle 180^\circ      & 10\angle 0^\circ \\[6pt]
15\angle 180^\circ      & 18.03\angle -33.69^\circ & 5\angle -90^\circ
\end{array}
\right]
$$

Para solucionarlo seleccionaremos el botón de cambiar tamaño y abrirá un menú [Fig \ref{fig:tamano}] donde escribiremos el valor de 2 siendo nuestra ecuación de 2 incógnitas y procederemos a darle en el botón ok.


Esto abrirá un nuevo menú con un total de 6 espacios disponibles, para ello podemos seleccionar casilla por casilla o usar la tecla $<\text{TAB}>$ para navegar entre casillas e ingresar nuestro sistema de ecuaciones, es importante mencionar que para los angulos se utilizo la c en lugar de $\angle$ por practicidad es decir un fasor de valor $5\angle45=5\text{c}45$.
Finalmente daremos en el botón resolver para observar los valores de las corrientes del circuito \ref{fig:Mesh} en la figura \ref{fig:PolarSOl}.

\begin{figure}[!h]
    \centering
    \includegraphics[width=0.9\linewidth]{Projects//Complex calc/SolPolares.png}
    \caption{Solución en forma polar}
    \label{fig:PolarSOl}
\end{figure}
\newpage
\subsection{Sistemas de Ecuaciones Reales}
Con el fin de demostrar su uso para calcular corriente de circuitos en corriente directa se armo el siguiente circuito [Fig. \ref{fig:DC}] considerando los valores del circuito pasado.
\begin{figure}[!h]
    \centering
    \includegraphics[width=0.6\linewidth]{Projects//Complex calc/DC.png}
    \caption{Circuito DC 2 Mallas}
    \label{fig:DC}
\end{figure}
$$
\begin{cases}
\begin{aligned}
(15 + 5)\,I_1 &\;-\; 15\,I_2 &=&\; 10 \\[6pt]
-15\,I_1 &\;+\; (15 + 10)\,I_2 &=&\; -5
\end{aligned}
\end{cases}
\Rightarrow
\left[
\begin{array}{cc|c}
20 & -15      & 10 \\[6pt]
-15     & 25 & -5
\end{array}
\right]
$$


Para solucionarlo seleccionaremos el botón de cambiar tamaño y abrirá un menú [Fig \ref{fig:tamano}] donde escribiremos el valor de 2 siendo nuestra ecuación de 2 incógnitas y procederemos a darle en el botón ok.

Esto abrirá un nuevo menú con un total de 6 espacios disponibles, para ello podemos seleccionar casilla por casilla o usar la tecla $<\text{TAB}>$ para navegar entre casillas e ingresar nuestro sistema de ecuaciones.
Finalmente daremos en el botón resolver para observar los valores de las corrientes del circuito \ref{fig:Mesh} en la figura \ref{fig:PolarSOl}.
\begin{figure}[!h]
    \centering
    \includegraphics[width=0.9\linewidth]{Projects//Complex calc/SolDC.png}
    \caption{Solución de la malla}
    \label{fig:PolarSOl}
\end{figure}

\section{Limitaciones}
Esta calculadora es capaz de resolver sistemas de ecuacuiones con números complejos asi como sistemas de ecuaciones lineales. Sin embargo, tiene la limitación de que sólo puede usarse para ese fin, no puede realizar ningún otro tipo de operación mas que resolver sistemas de ecuacion.

Otro limitación importante es el hecho de que sólo es posible realizar operaciones con sistemas de ecuaciones 10x10 como máximo, esto debido a que conlleva mucho tiempo de computo y requiere de gran espacio en pantalla.

\section{Posibles Errores}
Algunas de las posibles fallas al usar esta calculadora es el hecho d introducir erroneamenete algun valor.

Otro posible error es introduci sistemas sin solucion, para esta situacion el software lo indicará
\section{Licencia y acceso al código}

Este proyecto se distribuye bajo la \textbf{licencia MIT} \cite{mit_license}, lo que permite su uso, modificación y distribución libre, siempre que se mantenga el aviso de copyright.

El código fuente del proyecto está disponible públicamente en GitHub: 
https://github.com/DasReyxr/Py-ComplexCalc.

\begin{thebibliography}{9}
\bibitem{mit_license} 
MIT License. \url{https://opensource.org/licenses/MIT}

\bibitem{numpy} 
NumPy Libray. \url{https://numpy.org/}.

\end{thebibliography}


%%%%%%%%%%%%%%%%%%%%%%%%%%%%%%%%%%%%%%%%%%%%%%%%%%%%%%%%%%%%%%%%%%%%%%


\end{document}
